\section{Bounded positive height in spherical colors}
\label{sec:spherical}

We now prove the containment needed for a stabilizer inside a spherical
standard parabolic.  If the tests on the operated colors are not all
acyclic, the stabilizer lies in a conjugate of a proper standard
parabolic.  In this case, the main step rules out an infinite sequence
of positive words, each containing every operated color, along which
total test height is bounded.

Suppose such a sequence existed.  Positive amplitude gives a fixed
lower bound on generator levels; bounded total test height and the
fixed exterior generators give a fixed upper bound.  Every member of
the family therefore has its generator levels in one interval
$[d_1,d_2]$.  Meanwhile the greatest ordinary cochain degree of tested
cohomology drops by at least one after each word.  We will convert the
resulting test vanishing on a long cochain interval below all exterior
generators into an empty cochain placement.  The finite path bound
below removes one layer after another, at a uniformly bounded cost in
placement.  At a sufficiently late stage, nonzero test cohomology still
lies below the empty placement and the fixed $f$-generator lies above
it.  One transported projective would split, contradicting its
indecomposability.

\subsection{A bound on coefficient paths}

The underlying path calculation is the essential-path construction
associated with Ocneanu and the Jones--Wenzl projections; see
\cite[Sections~3.1--3.2]{Coquereaux2002} and
\cite[Theorem~5.4.3 and Proposition~5.6.10]{Cooper2007}.
We give its projection argument explicitly so that the resulting
nilpotence bound can be applied on an interval of cochain placements,
where the coefficient maps may depend on the placement.

\begin{lemma}\label{sph:path-bound}
Let \(Y\subseteq S\) be spherical.  There is an integer \(N=N(Y)\geq1\)
with the following property.  Consider a minimal bounded complex
of projectives and an integer interval of cochain placements on
which all its generator vertices lie in \(\widetilde Y\).
In each pure layer, every composition of \(N\) successive arrow
coefficient maps vanishes, provided all \(N+1\) placements of the
composition lie in that interval.
\end{lemma}

\begin{proof}
By Lemma~\ref{cat:unfolding}, the graph on \(\widetilde Y\)
is a disjoint union of finite simply laced graphs whose Cartan
matrices are positive definite.  We first establish a path
identity on one connected component with an edge.
Write \(\mathsf B\) for its adjacency matrix.

There is a strictly positive vector \(\mu\) with
\[
 \mathsf B\mu=\delta\mu,\qquad 0<\delta<2.
\]
For completeness, maximize the Rayleigh quotient of
\(\mathsf B\) on the unit sphere.  Replacing a maximizing
vector by its coordinatewise absolute value does not decrease
the quotient.  The resulting nonnegative eigenvector is
strictly positive: the eigenvalue equation at a zero coordinate
would force its neighbors to be zero, and connectedness would
then force the entire vector to vanish.  An edge makes the
maximum eigenvalue positive, and positivity of
\(2I-\mathsf B\) makes it less than \(2\).

Choose opposite arrow bases on each edge so that the trace
of their product, in either order, is
\(\sqrt{\mu_i\mu_j}\).  The nondegenerate opposite-arrow pairing
and ordinary symmetry of the Frobenius trace allow this choice.
For coefficient maps of a pure layer, the length-two component
of the square of the differential is therefore the relation
\begin{equation}\label{sph:eq:path-relation}
 \sum_{j\sim i}\sqrt{\mu_j/\mu_i}\,(i\,j\,i)=0.
\end{equation}
Here a path denotes the corresponding composition of coefficient
maps, at successive cochain placements.  Indeed, before division
by \(\mu_i\), its coefficient in the token trace is
\(\sqrt{\mu_i\mu_j}\).  Minimality excludes length-zero
differential entries, so only two length-one entries contribute
to this part of \(d^2\).  The relation is valid on the prescribed
interval because every intervening generator there is in
\(\widetilde Y\).

We prove that all sufficiently long abstract paths lie in the
span of insertions of \eqref{sph:eq:path-relation}.
Let \(\mathcal V_n\) be the finite-dimensional Hilbert space
with orthonormal basis the paths \((a_0,\ldots,a_n)\) of length
\(n\).  For \(n\geq2\) and \(1\leq i\leq n-1\), define
\[
 c_i^{(n)}:\mathcal V_n\longrightarrow\mathcal V_{n-2}
\]
by zero unless \(a_{i-1}=a_{i+1}\).  On a backtrack
\((a,b,a)\) at these positions, delete its two edges and
multiply the shortened path by \(\sqrt{\mu_b/\mu_a}\).
The adjoint inserts exactly the weighted sum in
\eqref{sph:eq:path-relation}, with the surrounding path fixed.
All these maps preserve the two endpoints.

The contraction and insertion operators satisfy
\begin{align}
 c_i^{(n)}(c_i^{(n)})^*
     &=\delta\,\mathrm{id}_{\mathcal V_{n-2}},
       \label{sph:eq:loop}\\
 c_i^{(n)}(c_{i-1}^{(n)})^*
     &=c_{i-1}^{(n)}(c_i^{(n)})^*
       =\mathrm{id}_{\mathcal V_{n-2}}
       \quad(2\leq i\leq n-1).
       \label{sph:eq:adjacent}
\end{align}
In \eqref{sph:eq:loop}, the scalar at the insertion vertex
\(a\) is \(\sum_{b\sim a}\mu_b/\mu_a=\delta\).
For \eqref{sph:eq:adjacent}, the subsequent contraction
forces the inserted detour to use the adjacent edge already
present in the path.  Its two factors are reciprocal square
roots, so their product is \(1\).  In particular, both
identities in \eqref{sph:eq:adjacent} use the same ambient
length \(n\).

Let
\[
 \mathcal E_n=\bigcap_{i=1}^{n-1}\ker c_i^{(n)},
 \qquad
 \mathcal E_0=\mathcal V_0,\quad
 \mathcal E_1=\mathcal V_1,
\]
and let \(\Pi_n\) be the orthogonal projection onto
\(\mathcal E_n\).  Since these spaces decompose by endpoints,
\(\Pi_n\) preserves endpoints.  We may therefore extend
\(\Pi_k\) to longer path spaces by applying it to the first
\(k\) edges while keeping the tail fixed.

Write
\[
 \delta=2\cos\theta,\qquad 0<\theta<\pi/2,\qquad
 [u]=\frac{\sin(u\theta)}{\sin\theta},
 \qquad h=\min\{j\in\mathbb Z_{\geq1}:j\theta\geq\pi\}.
\]
Then \([n]>0\) for \(1\leq n<h\), and
\([n]=\delta[n-1]-[n-2]\).
For \(2\leq n<h\), denote by \(\mathsf Q_n\) the extension
of \(\Pi_{n-1}\) to the first \(n-1\) edges of a length-\(n\)
path.  We claim
\begin{equation}\label{sph:eq:scalar}
 c_{n-1}^{(n)}\mathsf Q_n(c_{n-1}^{(n)})^*
     =\frac{[n]}{[n-1]}\Pi_{n-2}
     \quad\hbox{on }\mathcal V_{n-2},
\end{equation}
and
\begin{equation}\label{sph:eq:projection}
 \Pi_n=\mathsf Q_n-
       \frac{[n-1]}{[n]}\,
       \mathsf Q_n(c_{n-1}^{(n)})^*
                    c_{n-1}^{(n)}\mathsf Q_n.
\end{equation}

At \(n=2\), \(\mathsf Q_2\) is the identity, and
\eqref{sph:eq:scalar} is \eqref{sph:eq:loop}, since
\([2]/[1]=\delta\).
To pass from \(n-1\) to \(n\), let \(\mathsf R_n\) be the
extension of \(\Pi_{n-2}\) to the first \(n-2\) edges.
Extending \eqref{sph:eq:projection} for \(n-1\) by one edge
gives
\[
 \mathsf Q_n=\mathsf R_n-
       \frac{[n-2]}{[n-1]}\,
       \mathsf R_n(c_{n-2}^{(n)})^*
                    c_{n-2}^{(n)}\mathsf R_n.
\]
Endpoint preservation gives the identities
\[
 c_{n-1}^{(n)}\mathsf R_n
       =\Pi_{n-2}c_{n-1}^{(n)},\qquad
 \mathsf R_n(c_{n-1}^{(n)})^*
       =(c_{n-1}^{(n)})^*\Pi_{n-2}.
\]
They concern \(\mathsf R_n\), whose prefix stops before the
last pair, not \(\mathsf Q_n\).
Substitution, followed by \eqref{sph:eq:loop} and the two
adjacent identities, yields
\begin{align*}
 c_{n-1}^{(n)}\mathsf Q_n(c_{n-1}^{(n)})^*
 &=\delta\Pi_{n-2}
   -\frac{[n-2]}{[n-1]}\,
     \Pi_{n-2}c_{n-1}^{(n)}(c_{n-2}^{(n)})^*
                c_{n-2}^{(n)}(c_{n-1}^{(n)})^*\Pi_{n-2}\\
 &=\left(\delta-\frac{[n-2]}{[n-1]}\right)\Pi_{n-2}\\
 &=\frac{[n]}{[n-1]}\Pi_{n-2}.
\end{align*}
This proves \eqref{sph:eq:scalar}.

To justify the projection formula, set
\(T_n=c_{n-1}^{(n)}\mathsf Q_n\).
The identity just proved reads
\[
 T_nT_n^*=\frac{[n]}{[n-1]}\Pi_{n-2}.
\]
In finite-dimensional Hilbert spaces, an operator and its
product with its adjoint have the same image: their
orthogonal complements are both the kernel of the adjoint.
Thus \(T_n\) maps onto \(\mathcal E_{n-2}\).
Restricted to \(\operatorname{im}\mathsf Q_n\), its kernel
is \(\mathcal E_n\), because the prefix projection has
already imposed all contractions except the last one.
To verify the projection explicitly, put
\(\lambda=[n]/[n-1]>0\).
The established range identity gives
\[
 (T_n^*T_n)^2=\lambda T_n^*T_n,\qquad
 \mathsf Q_nT_n^*T_n=T_n^*T_n
       =T_n^*T_n\mathsf Q_n.
\]
Consequently
\(\mathsf Q_n-\lambda^{-1}T_n^*T_n\) is self-adjoint
and idempotent.  It has image in
\(\operatorname{im}\mathsf Q_n\cap\ker T_n\) and is
the identity on that subspace.  This is the projection
\eqref{sph:eq:projection}, completing the induction.

It remains to prove that $\mathcal E_n$ vanishes for some $n$.
The endpoint-dimension recurrence below will give that finite path bound.
Let \(\mathsf d_n(i,j)\) be the dimension of
\(\mathcal E_n\) with endpoints \(i,j\).
Extending a path of length \(n-1\) by one edge makes the
endpoint-dimension matrix of
\(\operatorname{im}\mathsf Q_n\) equal to
\(\mathsf d_{n-1}\mathsf B\).
The restricted surjection
\(T_n:\operatorname{im}\mathsf Q_n\to\mathcal E_{n-2}\)
preserves endpoints and has kernel \(\mathcal E_n\).
Hence, for \(2\leq n<h\),
\[
 \mathsf d_n=\mathsf d_{n-1}\mathsf B-\mathsf d_{n-2},
 \qquad \mathsf d_0=I,\quad\mathsf d_1=\mathsf B.
\]
Multiplication by \(\mu\) and induction give
\(\mathsf d_n\mu=[n+1]\mu\).
In particular, the recursion is valid at \(n=h-1\).
No denominator \([h]\) has been used.
Since \(\pi\leq h\theta<\pi+\theta<3\pi/2\),
\[
 0\leq\mathsf d_{h-1}\mu=[h]\mu\leq0.
\]
The entries of \(\mathsf d_{h-1}\) are nonnegative and
\(\mu\) is strictly positive.  It follows that
\(\mathsf d_{h-1}=0\), and therefore
\(\mathcal E_{h-1}=0\).  The same equality also forces
\([h]=0\); no classification of the graph was needed.

In finite-dimensional inner product spaces,
\[
 \mathcal E_n^\perp
       =\sum_{i=1}^{n-1}\operatorname{im}(c_i^{(n)})^*.
\]
Thus every individual path of length \(h-1\) is a linear
combination of insertions of
\eqref{sph:eq:path-relation}.  Evaluate such an identity
on the coefficient maps, starting at any fixed cochain
placement in the prescribed interval.  Each inserted
relation is then a valid length-two part of \(d^2=0\)
at the corresponding placement.  The coefficient maps
may depend on the placement; no translation invariance
of them is required.  The entire composition vanishes.

For a connected component without an edge, use the bound
\(N=1\).  There are finitely many components; take the
maximum of their bounds.  A longer path factors through
an initial segment of the relevant vanishing length,
and no path crosses components.  The empty set \(Y\)
is vacuous and may also be assigned \(N=1\).
\end{proof}

\subsection{Two height bounds force a gap}

Fix now a nonempty spherical \(Y\subseteq S\) and \(x\in A\).
Write \(H^{p,u}(e_\alpha E)\) for the cohomology in ordinary
cochain degree \(p\) and internal degree \(u\).
For \(a\in A_Y\), define the total test height
\begin{equation}\label{sph:eq:height}
 h(a)=\max\{p+u:H^{p,u}(e_\alpha E_L(ax))\ne0
                  \text{ for some }\alpha\in\widetilde Y,L\},
\end{equation}
with value \(-\infty\) if all these tests vanish.
There are finitely many labels and vertices, and all the
complexes are bounded, so this maximum is finite whenever
it is nonempty.
For the moment suppose that \(h(1)\) is finite.
Lemma~\ref{lay:height-change} shows that it remains finite
and nondecreasing under positive words in \(A_Y^+\).

\begin{lemma}\label{sph:cochain-drop}
Suppose that all tests on \(\widetilde Y\) in a bounded
family have ordinary cochain cohomology in degrees at
most \(m\).  Positive operations in \(Y\) preserve this
bound.  After a positive word containing every color of
\(Y\), the bound improves to \(m-1\).
\end{lemma}

\begin{proof}
For an acted vertex, Lemma~\ref{lay:positive} gives the
test shift \(\langle2\rangle[1]\), which lowers the
ordinary cochain bound by one.
At an unacted vertex \(\beta\), applying the vertex test
to the evaluation cone gives a triangle
\[
 \bigoplus_{\alpha\in\widetilde{\{i\}}}
      e_\beta P_\alpha\otimes_{\kk}e_\alpha E
       \longrightarrow e_\beta E
       \longrightarrow e_\beta F_iE .
\]
When \(\beta\) is adjacent to color \(i\), the first term
is a sum of acted tests shifted only in internal degree,
because each \(e_\beta P_\alpha\) is an arrow space.
It consequently has ordinary cochain cohomology at
most \(m\).  The long exact sequence gives zero above
\(m\) for the new test and a surjection
\[
 H^m_{\mathrm{coch}}(e_\beta E)
       \longrightarrow H^m_{\mathrm{coch}}(e_\beta F_iE),
\]
where each group retains its internal grading.
For a nonneighbor the first term vanishes and the test
is unchanged.

Thus once cohomology at degree \(m\) has vanished at a
vertex, later operations at other colors cannot restore
it: the later groups in that degree are quotients of
zero.  Another operation at its own color still gives
zero in degree \(m\), using the bound above \(m\).
After every color has been operated, degree \(m\) has
vanished everywhere.  The vertices of a single color
are pairwise nonadjacent, so the simultaneous color
operation satisfies the same conclusions.
\end{proof}

\begin{lemma}\label{sph:empty-placement}
Let a minimal bounded projective complex have generator
levels in \([d_1,d_2]\).  Suppose that on an integer
cochain interval \([l,r]\) all its generators lie in
\(\widetilde Y\), and all vertex tests on
\(\widetilde Y\) have zero cohomology in each ordinary
cochain degree in this interval, in every internal
degree.  With \(N=N(Y)\) from Lemma~\ref{sph:path-bound},
there are no generators at any placement
\[
 l+N(d_2-d_1+1)\leq p\leq r.
\]
\end{lemma}

\begin{proof}
Fix a vertex and a generator coefficient in layer \(d_2\)
at placement \(p\), with \(l<p\leq r\).
Multiplying its projective generator by its token gives
a cycle in the corresponding vertex test.  Every
differential entry is of positive path length and every
product of a token with such an entry is zero.
The tested cohomology vanishes at \(p\), so this token
coefficient is a boundary from placement \(p-1\).

Project such a boundary onto its token coefficient
space.  A length-one differential entry can contribute
only through an opposite arrow in the test, giving an
incoming arrow-coefficient image from the same layer.
A length-two entry can contribute only through a unit
path, and its source generator would lie in layer
\(d_2+1\).  This latter layer is absent.
These exhaust the possibilities, since minimality
excludes length-zero entries.  The perfect
opposite-arrow pairing therefore shows that the full
coefficient space at placement \(p\) is spanned by
incoming arrow-coefficient images at \(p-1\).

For \(p\geq l+N\), iterate this spanning assertion \(N\)
times.  Every path involved has its placements in
\([l,r]\), so Lemma~\ref{sph:path-bound} makes its
coefficient composition zero.  Layer \(d_2\) is
therefore empty for \(l+N\leq p\leq r\).

Proceed downward through the layers.  Suppose that
the higher layers have already been removed on
\([l+jN,r]\).  For layer \(d_2-j\) and
\(p>l+jN\), a possible token source from the next
higher layer at \(p-1\) is absent on that interval.
The same argument expresses the coefficient space
by incoming arrows, and then by \(N\)-step paths.
It removes this layer on \([l+(j+1)N,r]\).
The successive conclusions are
\[
 \begin{array}{c|c}
 \text{layer}&\text{interval proved empty}\\ \hline
 d_2 &[l+N,r]\\
 d_2-1 &[l+2N,r]\\
 \vdots&\vdots\\
 d_1 &[l+N(d_2-d_1+1),r].
 \end{array}
\]
On the last interval every possible generator layer
has been removed, proving the assertion.
\end{proof}

\begin{proposition}\label{sph:no-chain}
Assume \(h(1)>-\infty\).
There is no sequence
\[
 a_0=1,\qquad a_{k+1}=w_k a_k,\qquad
 w_k\in A_Y^+,\quad\operatorname{supp}(w_k)=Y,
\]
for which the integers \(h(a_k)\) have a common upper
bound.
\end{proposition}

\begin{proof}
Suppose that \(h(a_k)\leq h_*\) for all \(k\).
Choose a common lower bound \(d_1\) on the generator
levels of the initial family \(E_L(x)\).
Positive amplitude, Lemma~\ref{lay:positive} and the
bounded-layer formalism of Lemma~\ref{lay:structure},
preserves this bound at every stage.

All generator placements outside \(Y\) are fixed by
the family properties in Proposition~\ref{cat:family}.
Let \(d_{\mathrm{ext}}\) be their largest occupied
level.  The exterior includes the unchanged f
generator at placement \((0,0)\), so this is defined.
If a current global highest layer has degree
\(d>d_{\mathrm{ext}}\), it is supported entirely in
\(\widetilde Y\).  A nonzero such layer has a nonzero
top: incoming coefficient maps increase cochain
placement in a finite array and cannot generate a
nonzero array indefinitely.  Lemma~\ref{lay:extremal}
then gives a nonzero test in total degree \(d+2\).
It follows that every current generator has level
in the fixed interval
\[
 [d_1,d_2],\qquad
 d_2=\max(d_{\mathrm{ext}},h_*-2).
\]

By Lemma~\ref{sph:cochain-drop}, the common upper bound
on ordinary cochain test cohomology drops by at least
one at every step \(a_k\to a_{k+1}\).
It therefore tends to \(-\infty\).
Let \(p_{\mathrm{ext}}\) be the least exterior
generator cochain placement, over the initial finite
family.  It is fixed, and \(p_{\mathrm{ext}}\leq0\)
because of the f generator.
Set
\[
 N=N(Y),\quad
 K_0=N(d_2-d_1+1),\quad
 r=p_{\mathrm{ext}}-1,\quad l=r-K_0.
\]
At a sufficiently late stage all tested cohomology
on \(\widetilde Y\) lies in cochain degrees strictly
below \(l\).  Every generator in \([l,r]\) lies in
\(\widetilde Y\), since \(r<p_{\mathrm{ext}}\).
Lemma~\ref{sph:empty-placement} consequently shows
that every member of the family has no term at
cochain placement \(r=l+K_0\).

Some member nevertheless has nonzero test cohomology,
by positive height monotonicity from \(h(1)>-\infty\).
For that member there is a term in a cochain degree
below \(l\), while its unchanged f generator is in
degree \(0>r\).
The absent term in degree \(r\) splits its minimal
complex into two nonempty bounded minimal complexes,
one on either side of that degree.
Neither summand is contractible.  Indeed a contracting
homotopy would express its identity as \(dh+hd\);
every component of the latter lies in the additive
radical because all entries of \(d\) do.  Reduction
modulo the radical would give a nonzero identity
equal to zero.
This contradicts the indecomposability of
\(E_L(a_kx)\) from Proposition~\ref{cat:family}.
\end{proof}

\subsection{Positive extensions and a proper support}

\begin{lemma}\label{sph:common-extension}
For \(a,b\in A_Y\), there are \(p,q\in A_Y^+\) and
\(u\in A_Y\) such that
\[
 u=pa=qb,\qquad h(u)\leq\max(h(a),h(b)).
\]
The same upper bound holds for a common positive
left extension of any finite family.
\end{lemma}

\begin{proof}
The finite-type fraction property in
Theorem~\ref{cat:classical} first gives \(u=pa=qb\)
with \(p,q\) positive.  Whenever \(p\) and \(q\)
have a common left atom, write
\(p=\sigma_i p'\), \(q=\sigma_i q'\) and cancel it.
This replaces the common extension by
\(p'a=q'b\).
The sum of positive lengths decreases by two, so
the process terminates with no common left atom.

Suppose the resulting extension satisfies
\(h(u)>\max(h(a),h(b))\).
Choose one label \(L\) and one unfolded vertex
\(\alpha\in\widetilde Y\) attaining \(h(u)\) for
\(E_L(ux)\).
For this same family member, its starting test
height on each side is smaller than \(h(u)\).
If one starting set of tests vanished, persistence
of vanishing under positive words would contradict
the final nonzero test.  Otherwise
Lemma~\ref{lay:top-divisor} applies to each side.
It makes \(\alpha\) a left divisor of both unfolded
positive words.

Lemma~\ref{two:greedy} reflects this divisibility
to a common original color.  In its right-divisor
form this lemma is transferred to left divisors
by word reversal: every Artin braid relation is
preserved by reversal, and a color unfolds to a
product of commuting letters whose reversal
represents the same element.
Thus the original positive elements \(p,q\) have
a common left atom, a contradiction.
This proves the stated inequality.
Repeated pairwise application constructs a common
extension of any finite family, with height no
larger than the maximum of its starting heights.
\end{proof}

Positive support is well defined on \(A_Y^+\):
the positive monoid embeds, and every positive
Artin relation preserves the set of colors
occurring in a word.  In particular, the support
of a positive divisor is contained in the support
of the element it divides.

\begin{lemma}\label{sph:proper-support}
Assume \(h(1)>-\infty\), and put \(h_0=h(1)\).
There exist \(p\in A_Y^+\) with \(h(p)\leq h_0\)
and one proper subset \(J\subsetneq Y\) such that
\[
 b\in A_Y^+,\quad h(bp)\leq h_0
       \quad\Longrightarrow\quad
       \operatorname{supp}(b)\subseteq J.
\]
\end{lemma}

\begin{proof}
If no such \(p\) existed, then for every positive
\(p\) with \(h(p)\leq h_0\) the union of supports
of all its permitted continuations would equal
\(Y\).
For each color \(i\in Y\), choose a continuation
\(b_i\) whose support contains \(i\) and for which
\(h(b_i p)\leq h_0\).
Lemma~\ref{sph:common-extension} gives a common
positive left extension of these endpoints:
\[
 u=c_i b_i p\quad(i\in Y),\qquad
 c_i\in A_Y^+,\qquad h(u)\leq h_0.
\]
Cancelling \(p\) shows that \(w=up^{-1}=c_i b_i\)
is a positive element independent of \(i\).
Every \(b_i\) is a right divisor of \(w\), so
\(\operatorname{supp}(w)=Y\).
This constructs a new positive endpoint \(wp\)
with height at most \(h_0\).
At every intermediate letter of this operation
the height also stays at most \(h_0\), by positive
monotonicity.

Starting from \(p=1\) and repeating gives an
infinite sequence forbidden by
Proposition~\ref{sph:no-chain}.
Therefore some \(p\) has a proper union of the
supports of all its permitted continuations.
Take \(J\) to be precisely that union.
This establishes one common proper subset,
not merely a separately missing color for
each continuation.
\end{proof}

\begin{proposition}\label{sph:containment}
Let \(Y\subseteq S\) be spherical and \(x\in A\).
For the family stabilizer
\[
 H(Y,x)=\{s\in A_Y:E_L(sx)\simeq E_L(x)
                         \text{ for every }L\},
\]
one of the following holds:
\begin{enumerate}
 \item all tests on \(\widetilde Y\) at \(x\)
       are acyclic and \(H(Y,x)=A_Y\);
 \item there exist \(p\in A_Y^+\) and
       \(J\subsetneq Y\) such that
       \begin{equation}\label{sph:eq:containment}
        H(Y,x)\subseteq p^{-1}A_Jp .
       \end{equation}
\end{enumerate}
\end{proposition}

\begin{proof}
Suppose first that some test is nonzero.
Use \(p,J\) from Lemma~\ref{sph:proper-support}.
For \(s\in H(Y,x)\), the equivalences of family
members are preserved by every prefix of the
action.  In particular
\[
 h(ps)=h(p),\qquad h(bps)=h(bp)
                   \quad(b\in A_Y^+).
\]
Apply Lemma~\ref{sph:common-extension} to
\(p,ps\), obtaining
\[
 ap=bps,\qquad a,b\in A_Y^+,\qquad
 h(ap)=h(bps)\leq h(p)\leq h_0.
\]
Thus \(a\) is a permitted continuation of \(p\).
The equality \(h(bp)=h(bps)\) shows that \(b\)
is also permitted.  Both have support in \(J\),
so
\[
 s=p^{-1}b^{-1}ap\in p^{-1}A_Jp.
\]
This proves \eqref{sph:eq:containment}.

If all the tests are acyclic, the evaluation
term for each positive color twist is a bounded
acyclic complex of graded vector spaces tensored
with projectives.  It is contractible, since a
bounded acyclic vector-space complex splits
degreewise, in each internal degree.
The evaluation cone is consequently equivalent
to the original family member.
Every positive generator fixes the family, and
its inverse fixes it as well by applying the
inverse equivalence.
The generators therefore give
\(H(Y,x)=A_Y\).
This argument includes \(Y=\varnothing\), when
there are no tests and the group is trivial.
\end{proof}
